\uuid{8mP4}
\titre{Dérivées partielles, composée et structure d’une fonction }

\niveau{}
\module{ Fonctions de plusieurs variables }
\chapitre{Dérivation de fonctions composées}
\sousChapitre{Équations aux dérivées partielles}
\theme{}
\auteur{Maxime Nguyen}
\datecreate{ 2026-04-07}
\organisation{AMSCC}
\difficulte{3}

\contenu{

\texte{ 
Soit $f:\mathbb{R}^2\to\mathbb{R}$ une fonction de classe $C^1$ telle que
\[
\forall (x,y)\in\mathbb{R}^2,\qquad \frac{\partial f}{\partial x}(x,y)-\frac{\partial f}{\partial y}(x,y)=0.
\]

On introduit la fonction $h:\mathbb{R}^2\to\mathbb{R}$ définie par
\[
h(u,v)=f(u+v,u-v).
\]
}

\begin{enumerate}

\item   \question{Calculer $\dfrac{\partial h}{\partial v}(u,v)$ en fonction des dérivées partielles de $f$.}
\indication{}
\reponse{
On pose
\[
\varphi(u,v)=u+v,
\qquad
\psi(u,v)=u-v.
\]
Alors
\[
h(u,v)=f(\varphi(u,v),\psi(u,v)).
\]
Comme $f$ est de classe $C^1$, on peut appliquer la formule de dérivation des fonctions composées :
\[
\frac{\partial h}{\partial v}(u,v)
=
\frac{\partial f}{\partial x}(u+v,u-v)\frac{\partial \varphi}{\partial v}(u,v)
+
\frac{\partial f}{\partial y}(u+v,u-v)\frac{\partial \psi}{\partial v}(u,v).
\]
Or
\[
\frac{\partial \varphi}{\partial v}(u,v)=1,
\qquad
\frac{\partial \psi}{\partial v}(u,v)=-1.
\]
Donc
\[
\frac{\partial h}{\partial v}(u,v)
=
\frac{\partial f}{\partial x}(u+v,u-v)
-
\frac{\partial f}{\partial y}(u+v,u-v).
\]
}

\item   \question{En utilisant l’hypothèse sur $f$, montrer que
\[
\forall (u,v)\in\mathbb{R}^2,\qquad \frac{\partial h}{\partial v}(u,v)=0.
\]}
\indication{}
\reponse{
D’après la question précédente,
\[
\frac{\partial h}{\partial v}(u,v)
=
\frac{\partial f}{\partial x}(u+v,u-v)
-
\frac{\partial f}{\partial y}(u+v,u-v).
\]
Or l’hypothèse sur $f$ affirme que, pour tout couple $(x,y)\in\mathbb{R}^2$,
\[
\frac{\partial f}{\partial x}(x,y)-\frac{\partial f}{\partial y}(x,y)=0.
\]
En particulier, en prenant
\[
x=u+v,\qquad y=u-v,
\]
on obtient
\[
\frac{\partial h}{\partial v}(u,v)=0
\qquad\text{pour tout }(u,v)\in\mathbb{R}^2.
\]
}

\item   \question{{\textbf{Bonus:}} En déduire qu’il existe une fonction $g:\mathbb{R}\to\mathbb{R}$ telle que
\[
\forall (x,y)\in\mathbb{R}^2,\qquad f(x,y)=g(x+y).
\]}
\indication{}
\reponse{
Comme
\[
\frac{\partial h}{\partial v}(u,v)=0
\qquad\text{pour tout }(u,v)\in\mathbb{R}^2,
\]
la fonction $h$ ne dépend pas de la variable $v$. Il existe donc une fonction $g:\mathbb{R}\to\mathbb{R}$ telle que
\[
h(u,v)=g(u)
\qquad\text{pour tout }(u,v)\in\mathbb{R}^2.
\]

Or, par définition,
\[
h(u,v)=f(u+v,u-v).
\]
Pour relier cela à $f(x,y)$, on effectue le changement de variables
\[
x=u+v,\qquad y=u-v.
\]
En additionnant, on obtient
\[
x+y=2u,
\qquad\text{donc}\qquad
u=\frac{x+y}{2}.
\]
Ainsi
\[
f(x,y)=h\!\left(\frac{x+y}{2},\frac{x-y}{2}\right)=g\!\left(\frac{x+y}{2}\right).
\]

En posant alors
\[
k(t)=g\!\left(\frac{t}{2}\right),
\]
on obtient bien
\[
f(x,y)=g(x+y).
\]
Autrement dit, $f$ ne dépend que de la somme $x+y$.
}

\item   \question{Vérifier que toute fonction de la forme
$
f(x,y)=g(x+y)
$
avec $g\in C^1(\mathbb{R})$ vérifie bien
$$
\frac{\partial f}{\partial x}(x,y)-\frac{\partial f}{\partial y}(x,y)=0.
$$}
\indication{}
\reponse{
Supposons
\[
f(x,y)=g(x+y),
\]
où $k\in C^1(\mathbb{R})$.

Alors, en appliquant la formule de dérivation en chaîne,
\[
\frac{\partial f}{\partial x}(x,y)=g'(x+y)\cdot 1=g'(x+y),
\]
et
\[
\frac{\partial f}{\partial y}(x,y)=g'(x+y)\cdot 1=g'(x+y).
\]
On en déduit
\[
\frac{\partial f}{\partial x}(x,y)-\frac{\partial f}{\partial y}(x,y)
=
g'(x+y)-g'(x+y)=0.
\]
Ainsi, toute fonction de la forme
\[
f(x,y)=g(x+y)
\]
vérifie bien l’équation
\[
\frac{\partial f}{\partial x}(x,y)-\frac{\partial f}{\partial y}(x,y)=0.
\]
}

\end{enumerate}

}